% CardioTrace — main paper
% Last structural edit: 2026-04-18

\documentclass[journal,compsoc]{IEEEtran}

% ── Packages ──
\usepackage[utf8]{inputenc}
\usepackage[T1]{fontenc}
\usepackage{amsmath,amssymb}
\usepackage{graphicx}
\usepackage{booktabs}
\usepackage{multirow}
\usepackage{xcolor}
\usepackage{hyperref}
\usepackage{siunitx}
\usepackage{enumitem}
\usepackage{tabularx}
\usepackage[numbers]{natbib}

\newcommand{\system}{CardioTrace}
\newcommand{\ms}{\,\text{ms}}
\newcommand{\mv}{\,\text{mV}}
\newcommand{\bpm}{\,\text{bpm}}
\newcommand{\todo}[1]{\textcolor{red}{\textbf{[TODO: #1]}}}
\newcommand{\pending}[1]{\textcolor{orange}{\textbf{[PENDING: #1]}}}
\newcommand{\refplaceholder}[1]{\textcolor{blue}{[REF: #1]}}

\graphicspath{{figures/}}

% ============================================================
% PAPER ROADMAP — FUTURE WRITING NOTES (do not delete)
% ============================================================
% Source: 272-cell per-patient cardiologist read in
%   paper/headspace/12_systematic_misses/per_patient/
% Synthesis: paper/headspace/00_TOP_FINDINGS_v4_272cells.md
%
% The strongest pitch for this paper is the COMBINATION of:
%   (1) versioned prompts (v2, v5) + versioned reasoner patches
%       (v00..v10, with v11 as the targeted fix layer)
%   (2) the modular 5-layer architecture
%       (signal -> vision -> aggregator -> reasoner -> judge)
%   (3) the 272-cell manual failure-mode analysis that maps
%       every miss to a specific responsible layer
%
% This combination lets the paper claim something most ECG-LLM
% papers cannot: methodology for building diagnostic LLM
% pipelines with explicit per-layer failure-mode accountability.
%
% KEY POINTS TO BUILD OUT WHEN WRITING:
%
% (A) Per-layer failure-mode attribution table (paper figure)
%     Each of the 14 failure categories from the 272-cell read
%     maps to a specific layer. Suggested table columns:
%       Failure mode | Cell count | Responsible layer | Fix layer
%     Categories include:
%       - AFib false-positive on sinus-irregular (12+ cells)
%       - PVC architectural emission gap (20+ cells)
%       - Pacemaker detection gap (14+ cells)
%       - LBBB <-> RBBB <-> IVCD discrimination (10+ cells)
%       - Atrial flutter rate-window too narrow (7+ cells)
%       - Vision LM polarity inversion (suspect wider, 1 documented)
%       - Inferior MI under-commit (6+ cells)
%       - 2-deg / 3-deg AV block emission gap (4 cells)
%       - Ventricular aneurysm detection gap (3 cells)
%       - PTB-XL likelihood-0 phantom GT entries (30+ cells, GT-layer fix)
%       - LUDB chamber-enlargement over-coding (30+ cells, NOT our bug)
%       - LUDB axis-label disagreement (15+ cells, GT-quality)
%       - Chapman WPW/1AVB phantom labels (5+ cells, GT-quality)
%       - PAC/SVES architectural emission gap (6+ cells)
%
% (B) Prompt-evolution ablation (v2 vs v5 on same 485 patients)
%     Already in the data. Show how v5's "speaks to / out-of-scope"
%     prompt sections eliminated cross-talk noise but caused some
%     over-strictness on specific findings. Use ludb_143 as the
%     specific example (v5 catches LAE, v2 misses it). See
%     headspace/07_v2_vs_v5_prompts/README.md.
%
% (C) v11 as the targeted-fix layer
%     Phase 1 (highest impact, lowest risk):
%       - PTB-XL likelihood-0 phantom filter (GT loader, 3-line fix)
%       - PVC emission rule (vision + reasoner emission table)
%       - AFib false-positive override (vision-confirmed P-wave
%         in lead II overrides signal "sinus irregular")
%       - Vision-prompt polarity fix (above/below baseline instead
%         of upright/inverted)
%       - Pacemaker detection at signal pipeline
%     Phase 2 (medium): flutter window expansion, bundle-block
%       discrimination tuning, PAC/SVES emission, IMI threshold
%     Phase 3 (paper-novelty): 2-deg/3-deg AVB detection,
%       ventricular aneurysm rule, vision prompt v2 -> v5 alignment
%     Conservative v11 lift estimate: 50-80 captures across 970
%     cells (5-8% absolute improvement). Combined with Phase 1+2
%     pushes capture toward ~87-88% overall.
%
% (D) Honest GT-quality framing in Discussion
%     The headline capture rate (v2 81.5%, v5 79.4%) under-states
%     the model's real ECG-modality performance because some
%     "missed" findings are not detectable on the 12-lead ECG
%     alone (echo-derived LUDB chamber findings, PTB-XL phantom
%     likelihood-0 entries). After filtering phantoms and
%     accounting for multi-modal LUDB labels, effective
%     ECG-modality capture is approximately 87% (v2) / 85% (v5).
%
% (E) Suggested case-study figures
%     STRONG CATCHES (positive examples):
%       - ptbxl_2317 (Sokolow 3.87 mV definitive LVH)
%       - ptbxl_16748 (R/V5 3.17 mV massive LVH)
%       - ptbxl_12223 v2/v5 (textbook hyperacute T anterior STEMI)
%       - ptbxl_14534 v2/v5 (bifascicular block detection)
%       - ptbxl_11711 v2/v5 (WPW + AFib, clinically critical)
%     FAILURE CASE STUDIES (negative examples):
%       - ptbxl_15588 v2/v5 (5 false positives on a NORM patient,
%         AFib false-pos cascade)
%       - ptbxl_17719 v5 (LBBB missed despite clear broad-QS V1)
%       - ptbxl_12936 v2 (RBBB missed despite textbook rsR' M-shape)
%       - ludb_8 / ludb_93 (paced rhythm completely missed,
%         labeled "sinus regular" by signal pipeline)
%       - ptbxl_13499 v2/v5 (LV aneurysm pattern not committed,
%         classic persistent ST elevation + QS in V2)
%       - ludb_101 v5 (vision LM polarity inversion documented)
%
% (F) Suggested section structure that exploits all of the above
%     1. Architecture — describe 5 layers as independently
%        inspectable / patchable
%     2. Prompt evolution — v2 vs v5 head-to-head on same cohort
%     3. Per-layer failure-mode analysis — Table 1 from 272 cells
%     4. Targeted fixes (v11) — surgical fixes informed by failure
%        attribution; quantify capture lift per fix
%     5. v11 results — re-run on 970 cells; before/after capture
%        and which failure modes closed
%     6. Discussion — GT-quality framing for honest reporting
%
% ============================================================
% END PAPER ROADMAP NOTES
% ============================================================


% another thing we need to add to this paper we have to add the word deductive to showcasing the difference between the traditional approach and the new approach. and imply that old version was inductive.

\begin{document}

\title{\system{}: Traceable 12-Lead ECG Interpretation\\
via Vision-Language Fiducial Verification\\
and Retrieval-Augmented Diagnostic Reasoning}

\author{
  Amir~Sadjad~Taleban%
  \IEEEcompsocitemizethanks{
    \IEEEcompsocthanksitem Manuscript submitted April 2026.
    \IEEEcompsocthanksitem A.~S.~Taleban is with the Joseph J. Zilber School of Public Health, University of Wisconsin--Milwaukee, Milwaukee, WI 53201 USA.
  }
}

\IEEEtitleabstractindextext{%
\begin{abstract}
\todo{To be written last after all sections are finalized.}
\end{abstract}

\begin{IEEEkeywords}
electrocardiography, explainable AI, multimodal LLM, retrieval-augmented generation, fiducial detection, clinical decision support, open-source
\end{IEEEkeywords}
}

\maketitle
\IEEEdisplaynontitleabstractindextext

% ══════════════════════════════════════════════════════════════════════════
% INTRODUCTION
% ══════════════════════════════════════════════════════════════════════════
\section{Introduction}
\label{sec:intro}

\subsection{Clinical Importance of Accurate ECG Interpretation}

Electrocardiogram misinterpretation contributes to an estimated 4--12\% of adverse cardiac events in emergency settings, with missed ST-elevation myocardial infarction carrying a mortality rate three to four times higher than timely diagnosis~\cite{schlapfer2017computer}. Approximately 300 million ECGs are recorded annually worldwide, and studies consistently report that automated interpretation systems disagree with expert cardiologists in 20--30\% of cases, particularly for subtle findings such as early repolarization, borderline conduction abnormalities, and low-amplitude P-wave morphology~\cite{schlapfer2017computer}. The clinical stakes of ECG analysis --- where a single misclassified waveform can determine whether a patient receives emergent catheterization or is discharged --- motivate the development of systems that are not only accurate but transparently auditable.

\subsection{The Interpretability Gap in Automated ECG Analysis}

Automated 12-lead electrocardiogram interpretation has been a research objective since the 1960s, when the first computer-based ECG analysis programs were developed at Washington University and the Mayo Clinic~\cite{schlapfer2017computer}. Six decades later, the field has bifurcated. Commercial systems (GE Marquette 12SL, Philips DXL, Glasgow algorithm) rely on expert-authored decision logic that is transparent and clinically auditable but limited by the difficulty of encoding nuanced pattern recognition into fixed thresholds~\cite{schlapfer2017computer}. Deep learning classifiers~\cite{ribeiro2020automatic,hannun2019cardiologist,strodthoff2021deep} achieve higher aggregate accuracy on benchmark datasets but operate as end-to-end black boxes: a cardiologist presented with a neural network's ``atrial fibrillation'' label has no mechanism to inspect which features in which leads drove the decision, and no principled way to override a borderline call without repeating the analysis manually~\cite{rudin2019stop}. Recent work has quantified this trade-off directly: Patlatzoglou et~al.~\cite{explainability_cost_2025} demonstrated that imposing full explainability on ECG classifiers degrades performance by 13--43\%, suggesting that post-hoc interpretability methods exact a substantial accuracy cost.

This paper introduces \system{}, a system designed to sidestep this trade-off entirely. Rather than explaining a black-box classifier after the fact, \system{} constructs its diagnoses through an auditable chain of explicit measurements, vision-based morphological verification, and retrieval-augmented reasoning grounded in textbook criteria. The following subsections describe the landscape that motivates this design.

\subsection{Regulatory Pressure for Transparency}

This opacity is not merely an academic concern. The European AI Act classifies medical diagnostic AI as high-risk, requiring that users be provided with ``sufficient transparency to enable the deployer to interpret the system's output and use it appropriately''~\cite{eu_ai_act_2024}. The FDA's guidance on clinical decision support emphasizes that systems whose recommendations cannot be independently reviewed by clinicians require premarket clearance as medical devices~\cite{eu_ai_act_implications_2024}. Both regulatory frameworks create a practical incentive for ECG analysis systems that can explain their reasoning in terms a clinician can audit: which fiducial points were detected, which intervals were measured, which criteria for which conditions were met or not met, and what evidence supports each conclusion.

\subsection{The Necessity of Multimodal Morphological Reasoning}

A second limitation of current approaches, whether threshold-based or learned, is their single-modality design. Clinicians reading ECGs routinely combine signal-level analysis (measuring intervals, assessing axis) with visual morphological pattern recognition (the ``gestalt'' of a waveform shape, ST-segment curvature, T-wave symmetry) and encyclopedic disease knowledge (recalling that a bifid P-wave in lead~II with terminal negativity in V1 suggests left atrial enlargement). No existing automated system integrates all three reasoning modes~\cite{ansari2025ecg_llm_survey}. Threshold-based engines implement signal analysis but cannot see waveform shapes. Deep learning classifiers such as those benchmarked by Strodthoff et~al.~\cite{strodthoff2021deep} extract visual features implicitly but cannot articulate which morphological criteria drove a prediction. Neither consults a structured knowledge base during inference. The ECG-LLM survey by Ansari et~al.~\cite{ansari2025ecg_llm_survey} catalogues approximately 30 transformer- and LLM-based ECG methods; only one performs segment-level waveform decomposition, and none combine signal measurements, visual verification, and knowledge-grounded reasoning in a single pipeline.

\subsection{Multimodal LLMs as an Enabling Technology}

Recent advances in multimodal large language models have opened a new design space for ECG interpretation. PULSE~\cite{li2026pulse} fine-tunes a 7B vision-language model on over one million ECG image instructions, achieving 82.4\% AUC on PTB-XL but requiring large-scale supervised training and providing no mechanism for diagnostic traceability. ECG-LM~\cite{ecg_lm_2025} fuses raw ECG signals with a 7B language model via a ResNet encoder, achieving 0.693 zero-shot diagnostic accuracy but operating on whole-signal embeddings that obscure which morphological features drove each prediction. Q-HEART~\cite{qheart2025} combines an ECG encoder with FAISS-based retrieval of similar clinical reports, but its 1B model achieves only 0.329 accuracy on open-ended diagnostic queries. CKEPE~\cite{ckepe2024} uses GPT-4 to generate clinical-knowledge-enhanced prompts for zero-shot ECG classification (75.2\% AUC), but requires a proprietary cloud model at inference time.

A common limitation across these approaches is their treatment of ECG morphology as an implicit feature to be learned rather than an explicit measurement to be verified. None performs segment-level fiducial quality control, none retrieves structured diagnostic criteria from a curated knowledge base at inference time, and none produces diagnoses with citations traceable to specific textbook passages. Studies evaluating proprietary models directly on ECG images~\cite{seki2025ecg_vqa,lee2025ecg_mi,engelstein2025gpt4o_ecg} report that GPT-4o achieves only 30--66\% accuracy depending on task complexity, with the dominant failure mode being poor visual feature extraction from waveform morphology rather than insufficient medical knowledge. The emergence of open-source multimodal models in the 3B--27B parameter range --- including Gemma~3~\cite{gemma3_2025}, Qwen2.5-VL~\cite{qwen25vl2025}, and MiniCPM-V~\cite{minicpmv2024} --- creates an opportunity to deploy vision-language reasoning on clinic hardware without transmitting patient data externally, provided that prompts are engineered to compensate for these models' morphological limitations.

\subsection{Bridging Signal Processing and Visual Reasoning Through Domain Knowledge}

The morphological vocabulary of electrocardiography --- P-mitrale, RSR$'$, Wellens waves, Brugada pattern --- is drawn from a century of clinical observation codified in standard references including Goldberger's Clinical Electrocardiography, Chou's Electrocardiography in Clinical Practice, and Marriott's Practical Electrocardiography~\cite{goldberger2017clinical,marriott2019practical}. Prompt engineering for ECG interpretation must encode this domain-specific morphological language, not merely generic image descriptions. Studies evaluating generic ``describe this ECG'' prompts report accuracy below 60\% on fiducial verification tasks~\cite{li2026pulse,seki2025ecg_vqa}, whereas prompts that encode wave-specific expectations, lead-specific normal variants, and known algorithmic failure modes achieve 88--99\% accuracy (Section~\ref{sec:results_vision}). This observation motivates the segment-specific, context-aware prompt architecture described in Section~\ref{sec:vision}.
\subsection{Retrieval-Augmented Generation for Diagnostic Reasoning}

Retrieval-augmented generation (RAG) offers a mechanism for integrating structured disease knowledge into diagnostic reasoning without baking that knowledge into model weights. By retrieving relevant passages from a curated knowledge base at inference time, a RAG system can cite its sources, enabling post-hoc verification of the reasoning chain. Recent work has demonstrated the effectiveness of RAG in clinical settings: i-MedRAG~\cite{imedrag2025} improves medical QA accuracy by 3--7\% through iterative retrieval, while MedGraphRAG~\cite{medgraphrag2025} achieves 91.3\% on MedQA using hierarchical knowledge graphs with source-level citations. However, these systems operate on text-based clinical vignettes and multiple-choice benchmarks. To our knowledge, no prior work applies retrieval-augmented reasoning to ECG diagnostic interpretation, where the retrieved context must interface with signal-level measurements, vision model outputs, and structured morphological observations.

\subsection{Contributions}

This paper presents \system{}, a 12-lead ECG analysis system that integrates deterministic signal processing, multimodal vision-based quality control, and retrieval-augmented diagnostic reasoning into a single pipeline where every conclusion is fully traceable. The contributions are:

\begin{enumerate}[leftmargin=*]
\item \textbf{Vision-based fiducial verification}: To the best of our knowledge, the first application of multimodal LLMs as automated fiducial QC reviewers in electrocardiography, achieving 88--99\% agreement with expert annotations depending on the model.
\item \textbf{Segment-specific prompt architecture}: A wave-type-aware prompt design grounded in electrocardiographic textbook knowledge that transforms generic multimodal models into effective ECG morphology reviewers, with a 40-point accuracy improvement over generic prompts.
\item \textbf{Knowledge-grounded diagnostic reasoning}: A RAG pipeline backed by 94 structured disease documents derived from four standard ECG textbooks, producing diagnoses with explicit citations to specific criteria and measurements.
\item \textbf{Full diagnostic traceability}: An architecture in which every diagnostic statement traces from raw signal features through vision-verified morphology and retrieved reference passages to the stated conclusion, forming a complete audit chain amenable to clinical review.
\item \textbf{Multi-model, multi-dataset evaluation}: Comparison of five open-source vision models (3B--27B parameters) across PTB-XL, LUDB, and Chapman--Shaoxing, revealing that architectural design matters more than parameter count for ECG visual understanding.
\end{enumerate}


% ══════════════════════════════════════════════════════════════════════════
% METHODS
% ══════════════════════════════════════════════════════════════════════════
\section{Methods}
\label{sec:methods}

\system{} processes a standard 12-lead ECG through four stages designed to mirror clinical reasoning while maintaining full traceability. First, signal ingestion and preprocessing (\S\ref{sec:signal}) standardizes the raw recording and extracts a filtered signal for analysis alongside an unfiltered copy for amplitude measurements. Second, fiducial detection (\S\ref{sec:fiducial}) identifies P-wave, QRS, and T-wave boundaries via a consensus delineation algorithm, then segments each cardiac cycle into annotated strip images. Third, vision-based verification (\S\ref{sec:vision}) submits each segment image to a multimodal LLM with wave-specific prompts that encode domain knowledge from electrocardiographic textbooks, producing fiducial quality judgments and morphological annotations. Fourth, retrieval-augmented diagnostic reasoning (\S\ref{sec:rag}) assembles the vision outputs, extracted measurements, and beat-level narratives into a structured prompt, retrieves relevant passages from a 94-document disease knowledge base, and produces differential diagnoses with explicit citations to the retrieved criteria. The key design principle is that no stage operates as a black box: signal features are explicitly measured, vision verdicts are structured and auditable, and diagnostic conclusions cite specific knowledge-base passages.

\subsection{Signal Ingestion and Preprocessing}
\label{sec:signal}

Raw 12-lead ECG signals are loaded at their native sampling rate and converted to microvolts. When derived leads are absent, they are computed algebraically from leads~I and~II. A lead-name normalization layer handles database-specific conventions.

Signals are resampled to 500\,Hz via polyphase interpolation. Baseline wander is detected by computing the fractional power in the 0.05--0.5\,Hz band; when this exceeds 15\% of total signal power, a 600\ms{} median filter subtraction is applied. A 4th-order Butterworth bandpass filter (0.5--40\,Hz) is then applied with zero-phase (forward-backward) filtering. The outer 0.85\,s of each lead is excluded from analysis to avoid filtfilt edge transients, defining a safe analysis window. A second copy of the signal is retained at 500\,Hz without bandpass filtering (the ``morphology signal'') for amplitude measurements, ensuring that peak voltages are not attenuated by the filter.

\begin{figure*}[t]
\centering
\includegraphics[width=\textwidth]{fig1_preprocessing.png}
\caption{Preprocessing before and after: 12-lead ECG of patient~116 (PTB-XL). Left column shows the raw signal at native sampling rate; right column shows the preprocessed signal after resampling to 500\,Hz, baseline wander correction, and 0.5--40\,Hz bandpass filtering. Clinical ECG paper grid overlay at standard 25\,mm/s, 10\,mm/mV scale. Calibration bar (lower right): 200\,ms horizontal, 1\,mV vertical.}
\label{fig:preprocessing}
\end{figure*}

\subsection{Fiducial Detection and Cardiac Cycle Segmentation}
\label{sec:fiducial}

Fiducial detection uses the ecgdeli library's consensus delineation algorithm~\cite{pilia2021ecgdeli}, which performs independent P-wave, QRS, and T-wave delineation on each of 12 leads, then resolves disagreements through median-based consensus voting. The algorithm first establishes a consensus R-peak position across all leads, then validates P-wave and T-wave fiducials against this reference. The output is a 13-column fiducial point table (FPT) per lead per beat: P-onset, P-peak, P-offset, QRS-onset, Q-nadir, R-peak, S-nadir, QRS-offset (J-point), T-onset, T-peak, T-offset, and beat classification. Because algorithmic fiducial detection is known to be imprecise at wave boundaries --- particularly P-onset and T-offset, where inter-annotator variability is itself 10--30\ms{}~\cite{schlapfer2017computer} --- the vision-based verification stage (Section~\ref{sec:vision}) serves as an autonomous quality control layer that flags suspect placements before downstream measurements are computed.

Once R-peaks are established, each cardiac cycle is segmented into P, QRS, and T windows using percentage-based subdivision of the RR interval. The T-wave window, for example, begins at approximately 35\% of the RR interval after the R-peak and spans roughly 40\% of the RR interval, reflecting the known physiological proportions of the cardiac cycle~\cite{thygesen2018fourth}. Each segment is then rendered as an annotated strip image for submission to the vision model (Section~\ref{sec:vision}).
\begin{figure}[t]
\centering
\includegraphics[width=\columnwidth]{fig2_beat_segmentation.png}
\caption{Single cardiac cycle segmentation (lead~II, patient~116). Top: ECG waveform with P-wave (blue), QRS (red), and T-wave (green) windows highlighted, fiducial markers labeled. Bottom: percentage-based RR interval subdivision used to define each window.}
\label{fig:segmentation}
\end{figure}

\subsection{Vision-Based Verification and Morphological Annotation}
\label{sec:vision}

For each (segment, lead) pair, the cropped strip image is submitted to a small-to-medium-scale open-source multimodal LLM (3B--27B parameters) along with extracted metadata (fiducial times in ms, amplitudes, clinical context). The model returns a structured JSON response containing per-fiducial verdicts (correct, wrong with suggested correction, or uncertain), morphology observations (shape, polarity, symmetry), and clinical flags (ST elevation, T-wave inversion, pathological Q-wave depth). Table~\ref{tab:models} lists the five models evaluated.

\begin{table}[t]
\centering
\caption{Vision models evaluated for fiducial verification.}
\label{tab:models}
\begin{tabular}{llll}
\toprule
Model & Params & Architecture & Reference \\
\midrule
Qwen2.5-VL & 7B & ViT + MLP merger & \cite{qwen25vl2025} \\
Gemma-3 & 27B & SigLIP + local/global attn & \cite{gemma3_2025} \\
Gemma-3 & 12B & SigLIP + local/global attn & \cite{gemma3_2025} \\
MiniCPM-V & 3B & SigLIP + perceiver & \cite{minicpmv2024} \\
LLaVA & 13B & CLIP ViT-L + Vicuna & \cite{llavamed2023} \\
\bottomrule
\end{tabular}
\end{table}

\textbf{Prompt architecture.} Each wave type (P, QRS, T) receives a distinct prompt constructed from a four-layer hierarchical template: (1)~a \emph{domain context} layer encoding the morphological vocabulary and normal/variant expectations for that wave type, drawn from standard electrocardiographic textbooks~\cite{goldberger2017clinical,marriott2019practical}; (2)~a \emph{lead-specific expectations} layer encoding the expected polarity, amplitude range, and shape for the target lead; (3)~an \emph{algorithm explanation} layer describing how the upstream signal pipeline detected the fiducial and what failure modes are known, directing the vision model's attention to the points most likely to be incorrect; and (4)~a \emph{clinical context} layer injecting patient-level information such as detected conduction abnormalities (e.g., LBBB triggers supplementary context explaining that discordant ST-T changes are expected). This layered design is a critical architectural choice: generic ``analyze this ECG segment'' prompts produce accuracy below 60\% in our experiments, whereas the full prompt hierarchy achieves 88--99\% accuracy depending on the model (Section~\ref{sec:results_vision}).

\begin{figure*}[t]
\centering
\includegraphics[width=\textwidth]{fig3_vision_input.png}
\caption{Vision pipeline: segment-specific verification. Three cropped segment images (P-wave in lead~II, QRS in lead~V1, T-wave in lead~V3) as submitted to the Gemma-27B vision model, with abbreviated prompt excerpts (grey boxes) showing the four-layer prompt hierarchy and the model's structured JSON responses (yellow boxes) containing fiducial verdicts, morphological observations, and clinical flags.}
\label{fig:vision_input}
\end{figure*}

\textbf{Multi-model evaluation.} Five open-source multimodal models spanning 3B to 27B parameters were evaluated (Table~\ref{tab:models}). All models were deployed via the Ollama inference framework without external API calls. Gemma-27B and MiniCPM-V were deployed at scale across the full evaluation cohort; all five models were compared on a subset.

\subsection{Feature Extraction}
\label{sec:kb}

Approximately 200 clinical measurements are computed from the fiducial tables and signal amplitudes. Intervals (PR, QRS, QT) are taken as the median across all detected beats. QTc is computed using four correction formulas (Bazett, Fridericia, Framingham, Hodges). Amplitudes (R, S, T, Q) are measured from the unfiltered morphology signal to avoid bandpass attenuation. ST elevation and depression are measured at J+60\ms{} relative to a pre-QRS isoelectric baseline (median of the 100\ms{} segment preceding QRS-onset), using sex- and lead-specific thresholds consistent with AHA/ACC guidelines~\cite{thygesen2018fourth}: V1--V3 $\geq$0.2\mv{} (men) or $\geq$0.15\mv{} (women), all other leads $\geq$0.1\mv{}.

Morphological pattern classification assigns per-lead QRS shapes (QS, rS, RS, qR, RSR$'$, monophasic R, fragmented), ST curvature (concave, convex, horizontal), T-wave morphology (upright, inverted, biphasic, symmetric, asymmetric), and ST-T concordance. Voltage criteria for LVH are computed using Sokolow-Lyon, Cornell (sex-adjusted), and the Romhilt-Estes point score. Pathological Q-waves are identified by duration $\geq$40\ms{} or amplitude $\geq$25\% of R-wave, with the V1-V2 exception where any Q-wave is considered abnormal~\cite{thygesen2018fourth}. P-terminal force in V1 is computed as the product of the terminal negative deflection depth and duration (Morris criterion, threshold 0.04\,mV$\cdot$s).

\begin{figure}[t]
\centering
\includegraphics[width=\columnwidth]{fig4_kb_categories.png}
\caption{Knowledge base structure: 94 structured documents across seven diagnostic categories, with document counts and example conditions. Source material drawn from Goldberger, Chou, Marriott, and AHA/ACC guidelines.}
\label{fig:kb}
\end{figure}

\subsection{RAG-Based Diagnostic Reasoning}
\label{sec:rag}

The RAG layer receives the vision model's morphological annotations as its primary input, supplemented by the beat-level narrative and the extracted feature measurement set ($\sim$200 values). The vision output is central: the reasoner sees not only what the signal pipeline measured but how the vision model characterized each segment's morphology, including shape descriptions, polarity assessments, and clinical flags that the signal pipeline cannot produce.

\begin{figure*}[t]
\centering
\includegraphics[width=\textwidth]{fig5_rag_flow.png}
\caption{\system{} pipeline architecture. A 12-lead ECG is processed through signal preprocessing, fiducial detection, and beat segmentation (blue). Cropped segment images are verified by a vision LLM with segment-specific prompts (orange). A keyword-and-threshold retrieval mechanism selects relevant documents from the 94-document knowledge base (green). The reasoning LLM produces structured diagnoses with explicit citations to KB passages, measurements, and vision annotations (purple/gold).}
\label{fig:rag_flow}
\end{figure*}

\textbf{Knowledge base.} The disease knowledge base consists of 94 structured documents organized into seven categories: STAT/emergency conditions (14 files), rhythm disorders (15 files), conduction abnormalities (14 files), chamber enlargement (6 files), repolarization variants (5 files), ischemia/infarction patterns (15 files), and electrolyte/structural conditions (12 files), plus reference documents. Each document follows a standardized eight-section template: pathophysiology, electrical mechanism, lead-by-lead ECG manifestations, diagnostic criteria with numerical thresholds, reasoning complexity analysis, feature dependencies, retrieval requirements, and visualization specification. Source material was drawn from Goldberger's Clinical Electrocardiography, Chou's Electrocardiography in Clinical Practice, Marriott's Practical Electrocardiography, and the AHA/ACC/HRS guidelines~\cite{goldberger2017clinical,marriott2019practical,thygesen2018fourth}.

\textbf{Retrieval.} Document selection proceeds in two passes. First, keyword scanning of the narrative and vision annotations triggers condition-specific retrieval (e.g., ``ST elevated'' retrieves all STEMI variants, ``LBBB'' retrieves LBBB and Sgarbossa criteria). Second, a safety-net pass examines extracted features against threshold conditions (e.g., P-duration $\geq$120\ms{} triggers LAE documents) to capture conditions not explicitly mentioned in the narrative.

\textbf{Reasoning.} The selected documents, vision annotations, narrative, and measurements are assembled into a structured prompt submitted to a reasoning LLM. The architecture is model-agnostic; we evaluate DeepSeek-R1~\cite{deepseek_r1_2025} as the primary backend. The prompt enforces a systematic analysis protocol --- comprehensive beat survey followed by lead-by-lead analysis across all standard ECG parameters --- with embedded constraints that prevent common diagnostic pitfalls such as calling ischemia in the presence of expected discordant ST changes. The output is a structured JSON containing diagnoses with confidence levels, supporting evidence cited to specific leads and measurement values, criteria met and not met, borderline findings, excluded conditions with reasons, and an overall clinical interpretation. Each diagnosis links to the specific KB documents consulted, forming the traceability chain.

\subsection{Evaluation Datasets and Sampling Strategy}
\label{sec:datasets}

\textbf{PTB-XL.} The PTB-XL database~\cite{wagner2020ptbxl} contains 21,837 12-lead ECG records from 18,885 patients, with SCP-ECG diagnostic labels and the PTB-XL+ extension providing lead-level fiducial annotations. A stratified sample of 352 records was selected for evaluation, drawn to cover all target conditions with proportional disease representation. No model training or threshold tuning was performed on any partition.

\textbf{LUDB.} The Lobachevsky University Electrocardiography Database~\cite{kalyakulina2020ludb} provides 200 records with beat-level fiducial annotations across all 12 leads, making it the primary cross-dataset fiducial benchmark. A subset of 47 records received the full vision and RAG pipeline.

\textbf{Chapman--Shaoxing.} The Chapman--Shaoxing database~\cite{zheng2020chapman} contains 10,247 12-lead records with rhythm and morphology labels. A stratified sample of 87 records was selected to match the disease distribution of the PTB-XL evaluation set, ensuring that cross-dataset comparisons reflect detector generalization rather than prevalence differences.

\textbf{Sampling rationale.} The evaluation cohort of 486 patients (352 PTB-XL + 47 LUDB + 87 Chapman) was sized to provide statistically meaningful per-condition comparisons. For conditions with $\geq$12 true-positive cases, the Clopper-Pearson exact 95\% confidence interval on sensitivity has a margin of $\pm$22\% or better, which is standard for diagnostic pilot studies~\cite{flahault2005sample}. Table~\ref{tab:datasets} summarizes the evaluation cohort.

\begin{table}[t]
\centering
\caption{Evaluation dataset summary.}
\label{tab:datasets}
\begin{tabular}{lccc}
\toprule
Database & Total Records & Sampled & Unique Conditions \\
\midrule
PTB-XL & 21{,}837 & 352 & 16 \\
LUDB & 200 & 47 & 12 \\
Chapman--Shaoxing & 10{,}247 & 87 & 14 \\
\midrule
\textbf{Total} & \textbf{32{,}284} & \textbf{486} & --- \\
\bottomrule
\end{tabular}
\end{table}

\subsection{Performance Metrics}
\label{sec:metrics}

Performance is evaluated using sensitivity, specificity, positive predictive value (PPV), and F1 score, each reported with 95\% confidence intervals. For conditions with 12 or more true-positive cases in the evaluation cohort, confidence intervals are computed using the Clopper-Pearson exact method for proportions~\cite{newcombe1998ci} and the percentile bootstrap ($B = 2{,}000$) for F1. Conditions with fewer than 12 positive cases are reported as illustrative case series with full diagnostic reasoning chains rather than aggregate statistics, following STARD 2015 reporting guidelines~\cite{bossuyt2015stard}.

Vision model agreement across five architectures is compared using Cochran's Q test with post-hoc pairwise McNemar tests (Bonferroni-corrected). Cross-dataset heterogeneity is assessed using Fisher's exact test. Vision accuracy is computed as $\text{Accuracy}_{\text{vision}} = N_{\text{correct}} / N_{\text{total}}$, where a check is ``correct'' when the model's verdict matches the ground-truth fiducial placement. Results are stratified by segment type (P, QRS, T), database, and model.


% ══════════════════════════════════════════════════════════════════════════
% RESULTS
% ══════════════════════════════════════════════════════════════════════════
\section{Results}
\label{sec:results}

The evaluation spans three publicly available 12-lead ECG databases with distinct recording environments and annotation conventions (Table~\ref{tab:datasets}). PTB-XL (352 records) served as the primary benchmark; LUDB (47 records) and Chapman--Shaoxing (87 records) provided cross-dataset validation. All 486 records received the full pipeline including vision verification and RAG diagnostic reasoning. Detailed fiducial and interval measurement tables are reported in Appendix~A.

\subsection{Vision-Based Fiducial Verification}
\label{sec:results_vision}

Each segment is rendered as an annotated image and submitted to a multimodal LLM (3B--27B parameters) that judges whether the predicted fiducial markers are correctly positioned. To the best of our knowledge, this represents the first application of multimodal LLMs as automated fiducial quality-control reviewers in electrocardiography --- a task distinct from the ECG classification and report generation tasks evaluated in prior work~\cite{li2026pulse,ecg_lm_2025,ecg_chat2024,meit2024}.

\todo{Update numbers below once full 486-patient run completes.}
Across 3,352 segment-level checks on the PTB-XL cohort, the combined reviewer pool achieves 88.6\% agreement with expert-annotated fiducial positions. The accuracy hierarchy follows the expected morphological difficulty gradient: QRS segments at \textbf{96.0\%}, T-wave segments at \textbf{85.0\%}, and P-wave segments at \textbf{72.4\%}. The P-wave figure is notable: inter-annotator agreement among trained cardiologists for P-wave boundary placement ranges from 70\% to 85\%~\cite{schlapfer2017computer}, placing the automated reviewer within the human variability band.

Of the flagged checks, 11.4\% were marked as incorrectly placed, with segment-dependent rates: 27.6\% for P-wave, 15.0\% for T-wave, and 4.0\% for QRS. These correction rates provide a built-in quality metric absent from conventional fiducial pipelines, where errors propagate silently into downstream measurements.

Table~\ref{tab:vision} reports verification accuracy for the five models. Results are shown separately for PTB-XL and LUDB, which provide ground-truth fiducial annotations; Chapman--Shaoxing lacks fiducial labels and is excluded from this analysis.
\begin{table}[t]
\centering
\caption{Vision model fiducial verification accuracy.}
\label{tab:vision}
\begin{tabular}{lcccccc}
\toprule
Model & Params & P (\%) & QRS (\%) & T (\%) & Overall (\%) & N \\
\midrule
Qwen2.5-VL & 7B & --- & --- & --- & \textbf{99.3} & 549 \\
Gemma-27B & 27B & 71.6 & \textbf{99.1} & 88.6 & \textbf{91.5} & 2{,}803 \\
MiniCPM-V & 3B & 76.7 & 80.3 & 66.7 & 74.0 & 549 \\
Gemma-12B & 12B & --- & --- & --- & 68.7 & 549 \\
LLaVA-13B & 13B & --- & --- & --- & 45.2 & 527 \\
\bottomrule
\end{tabular}
\end{table}

Parameter count alone does not predict performance on this task: the 7B Qwen2.5-VL model outperforms every larger alternative, suggesting that architectural choices around fine-grained temporal signal tokenization matter more than raw capacity~\cite{li2026pulse,gem_ecg_2025}. The 46-point gap between the best and worst model (99.3\% vs 45.2\%) underscores that visual ECG understanding remains a differentiating benchmark for multimodal models.

Beyond fiducial verification, each model produces per-segment morphological annotations --- T-wave inversions, ST elevations, discordant ST-T patterns, and conduction morphology --- that feed directly into the RAG reasoning layer (Section~\ref{sec:results_rag}).

Beyond fiducial verification, the vision model produces per-segment morphological annotations: among the 2,731 segment assessments with Gemma flags, the model identified T-wave inversions, ST elevations, discordant ST-T patterns, and conduction morphology abnormalities. These qualitative findings feed directly into the RAG reasoning layer (Section~\ref{sec:results_rag}). Figure~\ref{fig:heatmap} illustrates a representative 12-lead heatmap for an inferior STEMI case (patient~7692, true positive), where abnormal and borderline flags distribute across leads and segments in an anatomically coherent pattern.

\begin{figure}[t]
\centering
\includegraphics[width=\columnwidth]{fig6_heatmap.png}
\caption{12-lead vision morphology assessment for patient~116 (Gemma-27B). Each lead shows three consecutive beats with colored overlays on P, QRS, and T segments: green = normal, yellow = borderline, red = abnormal, grey = not assessed. The ECG signal is plotted on a clinical grid background.}
\label{fig:heatmap}
\end{figure}

\subsection{Diagnostic Reasoning}
\label{sec:results_rag}

\system{} produces diagnostic labels through a retrieval-augmented generation pipeline that receives the vision model's morphological annotations as its primary input alongside extracted signal features and the beat-level narrative (see Section~\ref{sec:rag}). The reasoning LLM retrieves relevant passages from the 94-document knowledge base and synthesizes a differential diagnosis. On the PTB-XL cohort, the reasoner consulted a median of 20 documents per record and returned a median of 4 diagnoses (range 1 to 6).

Table~\ref{tab:rag} reports RAG diagnostic performance on PTB-XL alongside cross-dataset generalization on LUDB and Chapman--Shaoxing.

LVH, the condition with the most ground-truth positive cases in the pilot ($n = 17$), achieves sensitivity of 82.4\% and F1 of 0.65, demonstrating that the knowledge base effectively encodes the multi-criterion voltage assessment (Sokolow-Lyon, Cornell, Romhilt-Estes) that a single threshold cannot capture~\cite{thygesen2018fourth}. Inferior STEMI reaches 55.6\% sensitivity (F1 = 0.56), a plausible result given that ST-deviation narratives and reciprocal change patterns are well-represented in the KB. Anterior STEMI recall is 100\% but at the cost of catastrophic over-calling (45 FP, specificity 10.0\%), indicating that the KB's ST-elevation entries are over-represented. Deep learning classifiers on PTB-XL report macro-F1 of 0.70--0.80 across comparable condition sets~\cite{strodthoff2021deep,wagner2020ptbxl}; the RAG pipeline does not match that aggregate accuracy but produces a fully traceable reasoning chain (Section~\ref{sec:results_grounding}).

Cross-dataset transfer without threshold recalibration reveals predictable patterns. Rate-based detection is robust: sinus bradycardia achieves F1 of 0.96 on Chapman. Voltage-dependent conditions degrade across recording systems: LVH sensitivity drops to 11.1\% on LUDB and 45.7\% on Chapman, consistent with the well-documented sensitivity of voltage criteria to lead impedance and amplifier gain~\cite{schlapfer2017computer}.

\begin{table*}[t]
\centering
\caption{Diagnostic performance of the RAG pipeline on PTB-XL ($n = 352$) with cross-dataset generalization on LUDB and Chapman--Shaoxing. Sensitivity, specificity, and F1 are reported with Clopper-Pearson exact 95\% confidence intervals for conditions with $\geq$12 positive cases. Conditions with fewer positives are included for completeness but should be interpreted as case-level observations.}
\label{tab:rag}
\begin{tabular}{lccccccccc}
\toprule
Condition & GT+ & RAG Sens & RAG Spec & RAG F1 & LUDB Sens & LUDB F1 & Chapman Sens & Chapman F1 \\
\midrule
LVH & 17 & \textbf{0.824} & 0.727 & \textbf{0.651} & 0.111 & 0.183 & 0.457 & 0.356 \\
Inferior STEMI & 9 & \textbf{0.556} & 0.923 & \textbf{0.556} & --- & --- & --- & --- \\
Long QT & 4 & 0.500 & 0.789 & 0.222 & --- & --- & 0.286 & 0.229 \\
Lateral STEMI & 1 & 1.000 & 0.983 & 0.667 & --- & --- & --- & --- \\
AFib & 1 & 1.000 & 0.883 & 0.222 & 0.533 & 0.485 & 0.714 & 0.544 \\
RBBB & 3 & 0.333 & 0.983 & 0.400 & 0.250 & 0.200 & 0.500 & 0.510 \\
LAFB & 6 & 0.167 & 0.891 & 0.154 & --- & --- & --- & --- \\
Anterior STEMI & 11 & 1.000 & 0.100 & 0.328 & --- & --- & --- & --- \\
1st-deg AVB & 4 & 0.000 & 0.719 & 0.000 & --- & --- & 0.167 & 0.200 \\
WPW & 1 & 0.000 & 1.000 & 0.000 & --- & --- & --- & --- \\
Sinus bradycardia & --- & --- & --- & --- & 0.960 & 0.667 & \textbf{0.984} & \textbf{0.961} \\
LBBB & 0 & --- & 1.000 & --- & 0.500 & 0.174 & --- & --- \\
\bottomrule
\end{tabular}
\end{table*}

% Figures 7 and 8 removed — data is presented in Table~\ref{tab:rag}.

\subsection{Retrieval Grounding and Traceability}
\label{sec:results_grounding}

Every diagnosis produced by the RAG layer is accompanied by the specific knowledge-base documents the reasoner consulted, forming a complete provenance chain from extracted signal features through vision-verified morphology to retrieved reference passages and the stated diagnostic conclusion. Across the evaluation cohort, the median retrieval set is 20 documents (mean 19.5, max 25 of 94 available), with borderline and excluded conditions logged separately, providing a structured confidence gradient rather than a binary present/absent output. This per-diagnosis traceability constitutes a form of built-in interpretability that is structurally absent from end-to-end neural classifiers mapping directly from waveform to label~\cite{rudin2019stop}. Unlike post-hoc explanation methods such as Grad-CAM~\cite{ecgxplaim2025} or counterfactual generation~\cite{gcx2025}, which attempt to explain an opaque model's decision after the fact, \system{}'s traceability is intrinsic to the reasoning process itself --- the citations are not explanations of a black box but the actual evidence chain that produced the conclusion.

\subsection{Computational Performance}
\label{sec:results_timing}

Vision verification time scales linearly with the number of detected heartbeats. A typical ECG with 8--12 beats produces approximately 15 segment-lead checks per beat, totaling 120--180 individual vision calls per patient. Using Gemma-27B, each vision call averages 36 seconds, yielding 5--8 minutes per patient on a single GPU. The smaller Qwen2.5-VL (7B) completes checks substantially faster while achieving the highest accuracy (99.3\%). RAG reasoning adds 10--30 seconds per record depending on retrieval depth and reasoning complexity.

The full pipeline (vision + RAG) processes one patient in approximately 6--9 minutes on a single consumer GPU. These LLM-dependent stages are designed for asynchronous post-processing rather than real-time clinical use. Token consumption averages approximately 1,500 input tokens and 500 output tokens per vision check, and 8,000--15,000 input tokens per RAG query depending on the number of retrieved KB documents.

\todo{Add exact token usage statistics from the full 486-patient run.}


% ══════════════════════════════════════════════════════════════════════════
% DISCUSSION
% ══════════════════════════════════════════════════════════════════════════
\section{Discussion}
\label{sec:discussion}

\subsection{Related Work}

\system{} occupies a distinct position in the landscape of LLM-based ECG interpretation systems. Table~\ref{tab:comparison} compares the architectural properties of \system{} with the most closely related systems across six dimensions. We organize the comparison along three axes: modality and input representation, knowledge integration, and traceability.

\begin{table*}[t]
\centering
\caption{Architectural comparison of \system{} with related ECG-LLM systems. ``Trace.'' = full diagnostic traceability with citations to specific knowledge-base passages. ``Seg.'' = segment-level waveform analysis. ``VLM QC'' = vision-language model used for fiducial quality control.}
\label{tab:comparison}
\begin{tabular}{lccccccl}
\toprule
System & Input & Training & KB/RAG & Trace. & Seg. & VLM QC & Key Metric \\
\midrule
PULSE~\cite{li2026pulse} & ECG image & 1M+ inst. & None & No & No & No & 82.4\% AUC (PTB-XL) \\
Q-HEART~\cite{qheart2025} & Raw signal & 800K pairs & FAISS reports & Partial & No & No & 0.628 EM-Acc \\
CKEPE~\cite{ckepe2024} & Raw signal & 771K pairs & GPT-4 prompts & No & No & No & 75.2\% AUC (zero-shot) \\
ECG-LM~\cite{ecg_lm_2025} & Raw signal & Guideline-exp. & None & No & No & No & 0.693 Acc (PTB-XL) \\
ECG-Chat~\cite{ecg_chat2024} & Raw signal & 850K pairs & GraphRAG & Partial & No & No & 82.1\% faithfulness \\
MEIT~\cite{meit2024} & Raw signal & 800K reports & None & No & No & No & 0.610 BLEU-4 \\
GPT-4o~\cite{engelstein2025gpt4o_ecg} & ECG image & Zero-shot & None & No & No & No & 41--83\% (task-dep.) \\
\midrule
\textbf{\system{}} & \textbf{Seg. image} & \textbf{Zero-shot} & \textbf{94-doc KB} & \textbf{Yes} & \textbf{Yes} & \textbf{Yes} & \textbf{88--99\% (vision)} \\
\bottomrule
\end{tabular}
\end{table*}

\textbf{Input representation and training.} Most ECG-LLM systems operate on whole-signal embeddings (Q-HEART, CKEPE, ECG-LM, ECG-Chat, MEIT) or full ECG images (PULSE, GPT-4o), requiring large-scale supervised training to learn ECG representations. PULSE fine-tunes on over one million instruction samples; MEIT trains across 800K clinical reports. \system{} requires no ECG-specific training: it operates zero-shot on segment-level cropped images using open-source models (3B--27B), with domain knowledge encoded in prompts rather than model weights. This design enables immediate adaptation to new conditions by editing KB documents rather than retraining.

\textbf{Knowledge integration.} Q-HEART and ECG-Chat incorporate retrieval mechanisms, but retrieve similar clinical reports (Q-HEART) or use general-purpose GraphRAG (ECG-Chat) without structured diagnostic criteria. CKEPE uses GPT-4 to generate knowledge-enhanced prompts at test time, coupling the system to a proprietary cloud model. \system{}'s 94-document KB is structured around specific diagnostic criteria from four standard ECG textbooks, with each document containing actionable numerical thresholds. This structure enables single-pass keyword retrieval --- unlike iterative approaches such as i-MedRAG~\cite{imedrag2025} that require multiple retrieval rounds to compensate for imprecise initial queries --- because the KB's organization mirrors the clinical decision tree.

\textbf{Traceability.} No existing ECG-LLM system provides full diagnostic traceability. PULSE maps images to labels without intermediate reasoning. ECG-Chat achieves partial faithfulness through GraphRAG but cannot cite specific diagnostic criteria. GPT-4o studies report that the model produces ``fully accurate explanations'' in only 5\% of cases~\cite{lee2025ecg_mi}. \system{}'s architecture makes traceability structural rather than optional: each diagnosis links to specific KB document IDs, specific lead measurements, and specific vision model morphological observations. This property aligns with the agentic traceability demonstrated by DeepRare~\cite{deeprare2025} for rare diseases, adapted here to the structured domain of ECG interpretation.

\textbf{The explainability trade-off.} Patlatzoglou et~al.~\cite{explainability_cost_2025} demonstrated that imposing explainability constraints on ECG classifiers degrades performance by 13--43\% through information compression in learned representations. \system{} sidesteps this trade-off by constructing diagnoses from explicit measurements and retrieved criteria rather than compressing signals into latent spaces. The traceability cost is computational (minutes per patient vs.\ milliseconds for neural classifiers) rather than accuracy-based.

\subsection{RAG Diagnostic Reasoning: Strengths and Failure Modes}

The RAG pipeline demonstrates condition-dependent effectiveness that reflects the structure of the underlying knowledge base. LVH (F1 = 0.65) benefits from the KB's explicit encoding of multi-criterion voltage assessment. Inferior STEMI (F1 = 0.56) benefits from well-structured ST-deviation and reciprocal change patterns. The anterior STEMI over-calling (100\% sensitivity, 10\% specificity) reveals a retrieval calibration failure addressable through document weighting or retrieval score thresholds. LAFB (16.7\% sensitivity) and first-degree AVB (0\%) indicate insufficient KB coverage for these conditions.

Deep learning classifiers on PTB-XL report macro-F1 of 0.70--0.80~\cite{strodthoff2021deep,wagner2020ptbxl}. The RAG pipeline does not match that aggregate accuracy. The fundamental comparison, however, is not accuracy alone but the nature of the output: a neural classifier produces a label; \system{} produces a label, the evidence chain that led to it, the criteria that were and were not met, and the specific KB passages consulted. Whether this richer output justifies the accuracy gap is ultimately a clinical judgment.

\subsection{Vision Models as ECG Morphology Reviewers}

The 46-point inter-model accuracy spread (99.3\% to 45.2\%) across five models indicates that visual ECG understanding depends on architectural choices rather than parameter count alone. The 7B Qwen2.5-VL~\cite{qwen25vl2025} model's window attention mechanism and dynamic resolution processing appear well-suited to the fine-grained temporal features of ECG waveforms. The 27B Gemma-3~\cite{gemma3_2025} achieves strong per-segment results (99.1\% QRS) but lower overall accuracy due to P-wave difficulty (71.6\%). The 3B MiniCPM-V~\cite{minicpmv2024} offers a viable lightweight option at 74\% accuracy. LLaVA-13B's poor performance (45.2\%) is consistent with its general-purpose biomedical training, which does not include ECG-specific visual features.

The segment-specific prompt engineering is critical: without it, generic prompts produce accuracy below 60\%. This finding is consistent with the broader literature on medical prompt engineering, where domain-specific context dramatically improves performance~\cite{ckepe2024,seki2025ecg_vqa}.

\subsection{Cross-Dataset Generalization}

Rate-based detection transfers robustly across databases (sinus bradycardia F1 = 0.96 on Chapman). Voltage-dependent conditions degrade sharply: LVH sensitivity drops to 11.1\% on LUDB and 45.7\% on Chapman, consistent with the well-documented sensitivity of voltage criteria to lead impedance and amplifier gain~\cite{schlapfer2017computer}. This is a known limitation of all voltage-based ECG systems, not specific to \system{}.

\subsection{Limitations}

Several limitations should be noted. First, the evaluation cohort (486 patients) provides adequate statistical power for common conditions ($\geq$12 positives) but only case-level evidence for rare conditions such as WPW and lateral STEMI. Second, the reasoning LLM evaluation uses a single backend (DeepSeek-R1); comparison across multiple reasoning models is needed to assess architecture-agnosticism. Third, the knowledge base was authored by the research team from textbook sources rather than validated by an independent panel of cardiologists. Fourth, the vision and RAG stages add minutes of latency per patient, limiting the system to asynchronous analysis rather than point-of-care deployment. Fifth, the system has not been evaluated on pediatric, paced-rhythm, or artifact-heavy recordings.


% ══════════════════════════════════════════════════════════════════════════
% CONCLUSION
% ══════════════════════════════════════════════════════════════════════════
\section{Conclusion}
\label{sec:conclusion}

\system{} demonstrates that deterministic signal processing, multimodal vision-based quality control, and retrieval-augmented diagnostic reasoning can be combined into a single pipeline that produces fully traceable diagnostic conclusions. The vision layer provides autonomous fiducial verification at accuracy levels comparable to inter-annotator agreement among cardiologists (88--99\% depending on the model), and reveals that architectural design matters more than parameter count for ECG visual understanding. The RAG layer, backed by 94 structured documents derived from four standard ECG textbooks, produces diagnoses with explicit citations to specific criteria and measurements --- a form of built-in interpretability that is structurally absent from end-to-end classifiers.

The system's primary contribution is architectural: it establishes a pipeline design in which every diagnostic statement traces from raw signal features through vision-verified morphology and retrieved reference passages to the stated conclusion. This traceability property is independent of the specific models or thresholds used at each stage and provides a foundation for clinical audit, regulatory compliance, and iterative improvement. Among the five vision models evaluated, the 7B Qwen2.5-VL achieves the highest accuracy (99.3\%), demonstrating that effective ECG visual understanding is accessible to small, deployable models when paired with domain-specific prompt engineering.


% ══════════════════════════════════════════════════════════════════════════
% APPENDIX
% ══════════════════════════════════════════════════════════════════════════
\appendix

\section{Vision Model Input and Output Examples}
\label{app:vision}

\begin{figure*}[t]
\centering
\includegraphics[width=\textwidth]{fig_appendix_vision_segments.png}
\caption{Vision model input images and output by segment type. Each column shows one (patient, lead) pair. Row~1: P-wave segments. Row~2: QRS segments. Row~3: ST-T segments. Yellow boxes below each image contain the Gemma verdict (correct/wrong), morphology observations, and clinical flags.}
\label{fig:vision_segments}
\end{figure*}

\todo{Appendix Table~A1: Fiducial detection accuracy (PTB-XL fold~10, $n = 2{,}194$).}

\todo{Appendix Table~A2: Interval and amplitude measurement accuracy (PTB-XL fold~10).}


% ══════════════════════════════════════════════════════════════════════════
% REFERENCES
% ══════════════════════════════════════════════════════════════════════════
\bibliographystyle{IEEEtranN}
\bibliography{references}

\end{document}
